\section{Store sequences together with their metadata}
\label{sec:pg-storage-role}

A sequence search usually needs more than sequence symbols. A researcher may
search samples collected on a particular date, return their identifiers, and
join the hits to gene annotations. The database must keep these relationships
while allowing a search to read a small part of a long sequence.

Existing systems show why these relationships matter. BioWorkbench imports
variant and gene annotations into SQLite and joins them with
SQL~\autocite{Mondelli_2018}. The IBM Functional Genomics Platform uses a DB2
column store and mapping tables to connect genomes, genes, proteins, and
annotations~\autocite{seabolt2020ibmfunctionalgenomicsplatform}. Database-backed
sequence search has also been studied with DynamoDB records that pair each
sequence with an organism name~\autocite{flach2020nosql}. These examples motivate
keeping biological data queryable. Here the engineering question is how to
store the sequence itself so that database execution can request its symbols
as needed.

Consider a table \texttt{reads(id, sample\_id, seq)}. The following query
returns identifiers and does not need sequence symbols:

\begin{lstlisting}[style=bioseq-storage-sql,
  caption={A metadata query should not read sequence payloads.}]
SELECT id
FROM reads
WHERE sample_id BETWEEN 500 AND 600;
\end{lstlisting}

An alignment query on the same rows does need symbols, but it may request
several short intervals rather than the complete sequence. The sequence
interfaces in Section~\ref{sec:lazy-sequence-facade} already express those
requests. Storage must support them without first copying a whole chromosome
into one value.

PostgreSQL already stores large values out of line through TOAST. For some
uncompressed external values, a substring can fetch only the relevant
pieces~\autocite{postgresql17_toast}. BioSeqDB additionally needs packed
biological alphabets and a fragment reader usable by the alignment
algorithms. A sequence's packed payload can also exceed the ordinary
one-gigabyte limit for a PostgreSQL variable-length value. BioSeqDB therefore
separates three objects:

\begin{itemize}
  \item \textbf{The parent row} holds ordinary metadata and a small sequence
    locator. A locator identifies where the packed sequence is stored.
  \item \textbf{Fragment rows} hold fixed-size pieces of packed sequence
    data in internal PostgreSQL relations. A reader fetches fragment symbols
    only when they are needed.
  \item \textbf{Persistent indexes} keep lookup information derived from
    those symbols. Their identity and storage are described in the next
    chapter.
\end{itemize}

\warn{Adapted from Citus and Hydra Columnar}{
BioSeqDB's columnar layer adapts code from both Citus Columnar and Hydra
Columnar~\autocite{citus_columnar_storage,hydra_columnar_source}.
Hydra extended the Citus code with features including vectorized SQL
execution and a cache of decompressed column buffers. BioSeqDB modernized
the adapted C code into C++26 modules and substantially changed the storage,
reads, and updates. Sequence locators, fragment readers, and sequence
indexes are BioSeqDB-specific additions.}

The linked Hydra columnar documentation was unavailable when consulted for
this work. Understanding the chunk operations therefore required tracing
the available source code to reconstruct their write, read, and cache
lifetimes. The repository remains a public reference for the original
implementation~\autocite{hydra_columnar_source}.

The table access method connects this layout to PostgreSQL's table
operations~\autocite{postgresql17_tableam}. Citus illustrates the broader value of
building database features through PostgreSQL's extension interfaces,
including hooks and transaction callbacks~\autocite{citus_distributed_2021}.
The following sections explain the particular choices made in BioSeqDB.
